\BourbakiExerciseGroup

\begin{enumerate}
\def\labelenumi{\arabic{enumi})}
\tightlist
\item
  设 \(r_k\) 是整系数多项式，使得
\end{enumerate}

\[
p _ {k} = r _ {k} (s _ {1}, \dots , s _ {k}).
\]

\begin{enumerate}
\def\labelenumi{\alph{enumi})}
\item
  证明 \(r_k(\mathbf{Y}_1, \ldots, \mathbf{Y}_k)\) 包含项 \(\mathbf{Y}_1^k\) 与 \((-1)^{k-1} k\mathbf{Y}_k\) 。
\item
  证明：
\end{enumerate}

\[
r _ {k} \left(\mathrm{Y} _ {1}, \dots , \mathrm{Y} _ {k}\right) = \sum_ {p = 0} ^ {k - 1} (- 1) ^ {p} \quad \left| \begin{array}{c c c c c} \mathrm{Y} _ {p + 1} & \mathrm{Y} _ {p + 2} & \mathrm{Y} _ {p + 3} & \dots \mathrm{Y} _ {k - 1} & \mathrm{Y} _ {k} \\ 1 & \mathrm{Y} _ {1} & \mathrm{Y} _ {2} & \dots \mathrm{Y} _ {k - p - 2} & \mathrm{Y} _ {k - p - 1} \\ 0 & 1 & \mathrm{Y} _ {1} & \dots \mathrm{Y} _ {k - p - 3} & \mathrm{Y} _ {k - p - 2} \\ 0 & 0 & 0 & \dots \mathrm{Y} _ {1} & \mathrm{Y} _ {2} \\ 0 & 0 & 0 & \dots 1 & \mathrm{Y} _ {1} \end{array} \right|
\]

\begin{enumerate}
\def\labelenumi{\arabic{enumi})}
\setcounter{enumi}{1}
\item
  证明 \(n! s_n \equiv s_1^n\) 模 \(A[X_1, \ldots, X_n]\) 中由 \(X_1^2, X_2^2, \ldots, X_n^2\) 生成的理想。
\item
  设 \(f\) 是 \(\mathbf{K}[\mathbf{X}_1, \ldots, \mathbf{X}_n]\) 中的一个对称多项式，\(\varphi\) 是 \(\mathbf{K}[\mathbf{Y}_1, \ldots, \mathbf{Y}_n]\) 中使得 \(f(\mathbf{X}_1, \ldots, \mathbf{X}_n) = \varphi(s_1, \ldots, s_n)\) 的唯一多项式。证明 \(\varphi\) 的总次数等于 \(f\) 关于任意一个 \(\mathbf{X}_i\) 的次数。
\end{enumerate}

* 4) 设 K 是特征为 0 的域。证明如果 \(\alpha_{i}\) ( \(1 \leqslant i \leqslant n\) ) 是 K 中 \(n\) 个元素，使得 \(\alpha_{1}^{k} + \alpha_{2}^{k} + \cdots + \alpha_{n}^{k} = 0\) 对 \(n\) 个连续取值 \(h, h+1, \ldots, h+n-1\) 的 \(k\) 成立，则 \(\alpha_{1} = \alpha_{2} = \cdots = \alpha_{n} = 0\) 。用例子说明，如果 \(n\) 个 \(k\) 值不连续，或者 K 的特征不为 0，则该性质不再成立。

\begin{enumerate}
\def\labelenumi{\arabic{enumi})}
\setcounter{enumi}{4}
\tightlist
\item
  设 K 为一个交换域，且 \(f \in K[X_{1}, \ldots, X_{n}, Y_{1}, \ldots, Y_{n}]\) ；对每个置换 \(\sigma \in S_{n}\)，记 \(\sigma f\) 为有理分式
\end{enumerate}

\[
\sigma f \left(\mathrm{X} _ {1}, \dots , \mathrm{X} _ {n}, \mathrm{Y} _ {1}, \dots , \mathrm{Y} _ {n}\right) = f \left(\mathrm{X} _ {\sigma (1)}, \dots , \mathrm{X} _ {\sigma (n)}, \mathrm{Y} _ {\sigma (1)}, \dots , \mathrm{Y} _ {\sigma (n)}\right).
\]

则称 \(f\) 关于这些对 \((\mathbf{X}_i, \mathbf{Y}_i)\) 是对称的，如果 \(\sigma f = f\) 对所有 \(\sigma \in \mathfrak{S}_n\) 成立。

\begin{enumerate}
\def\labelenumi{\alph{enumi})}
\tightlist
\item
  对每个元素 \(\nu = (\lambda_1, \ldots, \lambda_n, \mu_1, \ldots, \mu_n)\)（属于 \(\mathbf{N}^{2n}\)）以及每个置换 \(\sigma \in \mathfrak{S}_n\)，记
\end{enumerate}

\[
\sigma^ {- 1} (\nu) = \left(\lambda_ {\sigma (1)}, \dots , \lambda_ {\sigma (n)}, \mu_ {\sigma (1)}, \dots , \mu_ {\sigma (n)}\right),
\]

使得 \(\mathfrak{S}_n\) 以这种方式成为 \(\mathbf{N}^{2n}\) 上的一个算子群。设 \(\gamma\) 为 \(\mathfrak{S}_n\) 在 \(\mathbf{N}^{2n}\) 中的任一轨道，并记 \(u_{\gamma}\) 为对称多项式

\[
\sum \mathrm{X} _ {1} ^ {\lambda_ {1}} \mathrm{X} _ {2} ^ {\lambda_ {2}} \dots \mathrm{X} _ {n} ^ {\lambda_ {n}} \mathrm{Y} _ {1} ^ {\mu_ {1}} \dots \mathrm{Y} _ {n} ^ {\mu_ {n}}
\]

，其中 \((\lambda_{1},\ldots,\lambda_{n},\mu_{1},\ldots,\mu_{n})\) 遍历 \(\gamma\) 。证明多项式 \(u_{\gamma}\) 构成关于 \((X_{i},Y_{i})\) 对称的所有多项式组成的向量空间在 K 上的一个基。

* b) 假设 K 的特征为 0。证明每个多项式 \(u_{\gamma}\) 都等于一个以对称函数为变量的、系数为有理数的多项式

\[
\boldsymbol {v} _ {\lambda \mu} = \sum_ {i = 1} ^ {n} \mathbf {X} _ {i} ^ {\lambda} \mathbf {Y} _ {i} ^ {\mu}
\]

（对 \((\lambda_i, \mu_i)\) 中不等于 \((0, 0)\) 者——在 \(\nu = (\lambda_1, \ldots, \lambda_n, \mu_1, \ldots, \mu_n)\) 中——的对数作归纳论证，并考虑乘积 \(u_\gamma v_{\lambda \mu}\)）。

\begin{enumerate}
\def\labelenumi{\alph{enumi})}
\setcounter{enumi}{2}
\tightlist
\item
  我们把 \((\mathbf{X}_i, \mathbf{Y}_i)\) 的初等对称函数理解为 \(n(n + 3)/2\) 个多项式 \(w_{hk} = u_\gamma\)，它们对应于轨道 \(\gamma\)——元素 \((\lambda_1, \ldots, \lambda_n, \mu_1, \ldots, \mu_n)\) 异于 \((0, \ldots, 0)\) 者——且满足：有 \(h\) 个对 \((\lambda_i, \mu_i)\) 等于 \((1, 0)\)，\(k\) 个等于 \((0, 1)\)，其余 \(n - h - k\) 个为 \((0, 0)\)。证明如果 \(\mathbf{K}\) 的特征为 0，则每个关于 \((\mathbf{X}_i, \mathbf{Y}_i)\) 对称的多项式都等于一个系数在 \(\mathbf{K}\) 中的、以初等对称函数为变量的多项式（考虑和式 \(\sum_{i=1}^{n} (\mathrm{UX}_i + \mathrm{VY}_i)^k\)，其中 \(\mathbf{U}\) 与 \(\mathbf{V}\) 是两个不定元，并利用 \(b\) )）。
\end{enumerate}

\(d)\) 如果 \(\mathbf{K}\) 是特征为 3 的域，且 \(n \geqslant 4\) ，证明 \(\sum_{i=1}^{n} \mathbf{X}_i^2 \mathbf{Y}_i^2\) 不等于对称函数 \(w_{hk}\) 的任何系数在 \(\mathbf{K}\) 中的多项式。

\begin{enumerate}
\def\labelenumi{\arabic{enumi})}
\setcounter{enumi}{5}
\tightlist
\item
  设 \(n\) 为一个整数 \(\geqslant 1\)，并设 \(X_1, \ldots, X_n\) 为不定元。对 \(1 \leqslant j \leqslant n\)，令 \(s_j\) 为次数为 \(j\) 的初等对称多项式（关于 \(X_i\)），且 \(u_j = (-1)^{j+1} s_j\) 。令 \(S(T) = \prod_{i=1}^{n} (1 - X_i T) = 1 - \sum_{j=1}^{n} u_j T^j\)，并对每个 \(k \geqslant 1\)，设 \(p_k = \sum_{i=1}^{n} X_i^k\) ；证明
\end{enumerate}

\[
p _ {k} = \sum a _ {\lambda_ {1} \lambda_ {2} \dots \lambda_ {n}} u _ {1} ^ {\lambda_ {1}} u _ {2} ^ {\lambda_ {2}} \dots u _ {n} ^ {\lambda_ {n}}
\]

求和遍及所有整数组 \((\lambda_{1},\ldots,\lambda_{n})\)（诸整数 \(\geqslant0\)）使得 \(\lambda_{1}+2\lambda_{2}+\cdots+n\lambda_{n}=k\)，且系数 \(a_{\lambda_{1}\lambda_{2}\ldots\lambda_{n}}\) 等于

\[
\frac {k \left(\lambda_ {1} + \lambda_ {2} + \cdots + \lambda_ {n} - 1\right) !}{\lambda_ {1} ! \lambda_ {2} ! \dots \lambda_ {n} !}
\]

（把级数 \(S'(T) / S(T)\) 展开为 \(T\) 的形式幂级数）。

\begin{enumerate}
\def\labelenumi{\arabic{enumi})}
\setcounter{enumi}{6}
\tightlist
\item
  沿用练习 6 的记号，证明
\end{enumerate}

\[
\sum a _ {\lambda_ {1} \lambda_ {2} \dots \lambda_ {n}} = - 1 + k \sum_ {j = 0} ^ {k / (n + 1)} \frac {(- 1) ^ {j}}{k - j n} \binom {k - j n} {j} 2 ^ {k - j (n + 1)}.
\]

\begin{enumerate}
\def\labelenumi{\arabic{enumi})}
\setcounter{enumi}{7}
\tightlist
\item
  令
\end{enumerate}

\[
s _ {k} = \sum \mathrm{A} _ {m _ {1} \dots m _ {n}} p _ {1} ^ {m _ {1}} \dots p _ {n} ^ {m _ {n}},
\]

其中 \(\mathbf{A}_{m_1\dots m_n}\) 是有理数，且求和遍及所有族 \((m_1,\dots,m_n)\) ，它们满足 \(\sum i\cdot m_i = k\)。证明 \(\sum \mathbf{A}_{m_1\dots m_n} = 0\) 对 \(k > 1\) 成立，且 \(\sum |\mathbf{A}_{m_1\dots m_n}| = 1\) 对 \(k\geqslant 1\) 成立。由此推出，在写出 \(q_{j} = -p_{j}\) 后我们有

\[
(- 1) ^ {k} s _ {k} = \sum \left| \mathrm{A} _ {m _ {1} \dots m _ {n}} \right| q _ {1} ^ {m _ {1}} \dots q _ {n} ^ {m _ {n}}.
\]

\begin{enumerate}
\def\labelenumi{\arabic{enumi})}
\setcounter{enumi}{8}
\tightlist
\item
  \begin{enumerate}
  \def\labelenumii{\alph{enumii})}
  \tightlist
  \item
    设 \(\mathbf{K}\) 为一个域，\(\mathbf{X}_1, \ldots, \mathbf{X}_n\) 为不定元。证明对任意指数 \(\nu_1 < \nu_2 < \ldots < \nu_n\)，行列式 \(\det(\mathbf{X}_i^{\nu_j})\) 在 \(\mathbf{K}(\mathbf{X}_1, \ldots, \mathbf{X}_n)\) 中非零。
  \end{enumerate}
\end{enumerate}

\begin{enumerate}
\def\labelenumi{\alph{enumi})}
\setcounter{enumi}{1}
\tightlist
\item
  设 \(s_k\) ( \(1 \leqslant k \leqslant n\) ) 为 \(X_j\) 中的初等对称多项式，且 \(r_1, \ldots, r_n\) 为 \(n\) 个元素，均属域 \(K(X_1, \ldots, X_n)\) 。我们令（引入新的不定元 T）
\end{enumerate}

\[
\mathrm{R} (\mathrm{T}) = \frac {1 + r _ {1} \mathrm{T} + \cdots + r _ {n} \mathrm{T} ^ {n}}{1 + s _ {1} \mathrm{T} + \cdots + s _ {n} \mathrm{T} ^ {n}} = 1 + u _ {1} \mathrm{T} + \dots + u _ {k} \mathrm{T} ^ {k} + \dots .
\]

证明如果有 \(n\) 个系数 \(u_{k}\) 为零，则所有 \(u_{k}\) 都为零。（把 \(R(T)\) 分解为部分分式。）

* \(c\) ) 假设 \(\mathbf{K}\) 的特征为零。令 \(\{\nu_1,\nu_2,\dots,\nu_n\}\) 为 \(n\) 个 \(\geqslant 1\) 的 \(\mathbf{N}\) 中元素构成的集合，并设 \(\mathbf{M}\) 为该集合在 \(\mathbf{N}\) 中的补集；再假设 \(\mathbf{M}\) 在加法下稳定。证明若 \(y_{1},\ldots ,y_{n}\) 是 \(n\) 个元素，均属 \(\mathbf{K}(\mathbf{X}_1,\dots,\mathbf{X}_n)\) ，使得

\[
p _ {\nu _ {j}} (y _ {1}, \dots , y _ {n}) = p _ {\nu _ {j}} (\mathbf {X} _ {1}, \dots , \mathbf {X} _ {n})
\]

对 \(1 \leqslant j \leqslant n\) 成立，则 \(y_{j} = X_{j}\) 对 \(1 \leqslant j \leqslant n\) 成立，至多相差一个置换。（考察 \(\log R(T) = v_{1}T + v_{2}T^{2} + \cdots + v_{k}T^{k} + \cdots\) 的展开式，注意 \(v_{\nu_{j}} = 0\) 对 \(1 \leqslant j \leqslant n\) 成立，并考察 \(v_{k}\) 与 \(u_{k}\) 之间的关系。） *

\begin{enumerate}
\def\labelenumi{\arabic{enumi})}
\setcounter{enumi}{9}
\item
  设 \(a \in \mathbf{A}\)。对每个多项式 \(h \in \mathbf{A}[\mathbf{X}]\)，令 \(h_a(\mathbf{X}) = h(\mathbf{X} + a)\) ；证明 \(\operatorname{res}(f_a, g_a) = \operatorname{res}(f, g)\)，\(\operatorname{dis}(f_a) = \operatorname{dis}(f)\) 。
\item
  在本练习中，我们假设 A 是一个域， \(f(0) \neq 0\) 且 \(\deg g < \deg f = p\) 。设 \(\sum \alpha_{n} X^{n}\) 为有理分式 \(g(X) / f(X)\) 的形式幂级数展开；对每个整数 \(k \geqslant 0\)，令 \(H_{0}^{(k)}\) 为 IV 第 90 页练习 1 中定义的汉克尔行列式。a) 证明 \(H_{0}^{(k)} = 0\) 对 \(k > p\) 成立，并把 \(H_{0}^{(p)}\) 作为 \(\operatorname{res}(f, g)\) 的函数计算出来。
\end{enumerate}

\begin{enumerate}
\def\labelenumi{\alph{enumi})}
\setcounter{enumi}{1}
\tightlist
\item
  在什么条件下我们有 \(\mathbf{H}_0^{(p)} = \mathbf{H}_0^{(p - 1)} = \ldots = \mathbf{H}_0^{(p - r + 1)} = 0\) ？
\end{enumerate}

\begin{enumerate}
\def\labelenumi{\arabic{enumi})}
\setcounter{enumi}{11}
\tightlist
\item
  存在唯一的多项式 \(P_{p,q}(T_{0},\ldots,T_{p},U_{0},\ldots,U_{q})\)，它具有整数系数，使得对每个环 A 以及任意多项式 \(f=t_{p}X^{p}+\cdots+t_{0}, g=u_{q}X^{q}+\cdots+u_{0}\)（属于 A{[}X{]}），我们有
\end{enumerate}

\[
\operatorname{res} _ {p, q} (f, g) = \mathrm{P} _ {p, q} \left(t _ {0}, \dots , t _ {p}, u _ {0}, \dots , u _ {q}\right).
\]

如果我们把权 i 赋予 \(T_{i}\) 与 \(U_{i}\)，则 \(P_{p,q}\) 是权为 pq 的等权式。

\begin{enumerate}
\def\labelenumi{\arabic{enumi})}
\setcounter{enumi}{12}
\tightlist
\item
  设 \(n \geqslant 1\) 是一个整数，并将 \(\mathbf{N}^n\) 的元素按字典序排列（记作 \(\alpha \preccurlyeq \beta\) ）。对每个 \(\alpha \in \mathbf{N}^n\) ，令 \(T_{\alpha}\) 为 \(\mathbf{Z}[X_1, \ldots, X_n]\) 中形如 \(X^{\alpha} + \sum_{\beta < \alpha} u_{\alpha\beta} X^{\beta}\) 的元素组成的集合（其中 \(u_{\alpha\beta} \in Z\) ).
\end{enumerate}

\begin{enumerate}
\def\labelenumi{\alph{enumi})}
\tightlist
\item
  证明 \(\mathbf{T}_{\alpha}\) 的一个元素与 \(\mathbf{T}_{\beta}\) 的一个元素的乘积属于 \(\mathbf{T}_{\alpha + \beta}\) 。\\
\item
  设 \(\varepsilon_{k}\) 为元素 \((1, 1, \ldots, 1, \underbrace{0, \ldots, 0}_{k})\)（属于 \(\mathbf{N}^{n}\)）。证明 \(s_{k} \in \mathbf{T}_{\varepsilon_{k}}\) 。由此推出
\end{enumerate}

\(S(\alpha) = \prod_{k=1}^{n} s_k^{\alpha_k - \alpha_{k+1}} \text{ 属于 } T_\alpha, \text{ 其中 } \alpha \in \mathbb{N}^n \text{ 满足 } \alpha_1 \geqslant \alpha_2 \geqslant \ldots \geqslant \alpha_n. (\text{ 令 } \alpha_{n+1} = 0.)\) \(c) \text{ 沿用 IV 第 66 页的记号，推出 } c_{\alpha\alpha} = d_{\alpha\alpha} = 1 \text{ 对 } \alpha \in D_k \text{ 且 } c_{\alpha\beta} = d_{\alpha\beta} = 0 \text{ 若 } \alpha < \beta.\)

\begin{enumerate}
\def\labelenumi{\arabic{enumi})}
\setcounter{enumi}{13}
\tightlist
\item
  设 \(L\) 是一个有限有序集，\(A\) 是一个交换环。在环 \(M\) （由 \(A\) 上的方阵组成，以 \(L\) 为指标集）中，矩阵 \(m = (m_{uv})\) （满足 \(m_{uv} = 0\)，除非 \(u \leqslant v\)）的集合是一个子环 \(B\) 。
\end{enumerate}

\begin{enumerate}
\def\labelenumi{\alph{enumi})}
\item
  元素 \(m = (m_{uv})\)（属于 B）在 B 中可逆，当且仅当 \(m_{uu}\) 对所有 \(u \in L\) 在 A 中可逆。
\item
  今后假设 \(\mathbf{A} = \mathbf{Z}\)。设 \(\zeta\) 为 \(\mathbf{B}\) 中由 \(\zeta(u, v) = 1\)（若 \(u \leqslant v\)）且 \(\zeta(u, v) = 0\)（否则）定义的元素。证明 \(\zeta\) 有一个逆 \(\mu\)（也写作 \(\mu_{\mathrm{L}}\)），位于 \(\mathbf{B}\) 中（称为 L 的默比乌斯函数）。设 \(f\) 与 \(g\) 是从 L 到 \(\mathbf{Z}\) 中的两个映射。证明关系式 \(\ll f(x) = \sum_{y \geqslant x} g(y)\)（对所有 \(x \in L\)）与 \(\ll g(x) = \sum_{y \geqslant x} \mu(x, y) f(y)\)（对所有 \(x \in L\)）是等价的。
\item
  假设 \(L\) 是有限个有序集 \(L_1, \ldots, L_n\) 的乘积。证明 \(\mu_L(x, y) = \prod_{i=1}^{n} \mu_{L_i}(x_i, y_i)\) 对 \(x = (x_1, \ldots, x_n)\) 与 \(y = (y_1, \ldots, y_n)\)（满足 \(x \leqslant y\)）成立。
\item
  设 \(m \geqslant 1\) 是一个整数，并令 \(D\) 为 \(m\) 的因子组成的集合，配备序关系 « \(x\) 整除 \(y\) 」。证明 \(\mu_{D}(x, y) = (-1)^{s}\)，其中 \(s\) 是素数 \(p\) （整除 \(y/x\) 者）组成的集合的基数（当 \(y/x\) 不被任何整数的平方整除时），而否则 \(\mu_{D}(x, y) = 0\)。（通过应用 \(c\) ）化简为 \(m\) 是一个素数的幂的情形。）
\item
  假设有序集 L 是格序的。设 \(a, b, c\) 是 L 中互不相同的元素，使得 \(a \leqslant b \leqslant c\)，并设 \(\Phi\) 为所有 \(x \in \mathbf{L}\) （满足 \(a \leqslant x\) 且 \(\sup(x, b) = c\)）组成的集合，\(\Psi\) 为相应的由 \(y \leqslant c\) 且 \(\inf(y, b) = a\) 定义的集合。证明我们有关系式 \(\sum_{x \in \Phi} \mu(a, x) = \sum_{y \in \Psi} \mu(y, c) = 0\) 。
\item
  设 T 是一个有限集合，\(\mathcal{R}(T)\) 是 T 上等价关系的集合，按关系 « R 比 S 更细 » 排序，记作 \(R \leqslant S\) 。证明 \(\mathcal{R}(T)\) 是格序的。设 \(\mathcal{R}\) 中的 R, S 满足 \(R \leqslant S\) ；设 \(U_1, \ldots, U_k\) 为 S 下的等价类，且对 \(1 \leqslant i \leqslant k\)，设 \(\alpha_i\) 为包含在 \(U_i\) 中的 R 下等价类的数目。证明 \(\mu_{\mathcal{R}(T)}(R, S) = \prod_{i=1}^{k} (-1)^{\alpha_i - 1} (\alpha_i - 1)!\)（若 \(A_i\) 是 R 下的等价类组成的集合
\end{enumerate}

包含在 \(U_{i}\) 中，观察到区间 \([R, S]\) （在 \(\mathcal{R}(T)\) 中）同构于 \(\mathcal{R}(A_{1}) \times \ldots \times \mathcal{R}(A_{k})\)，应用 c) 并化简到 R 是 \(\mathcal{R}(T)\) 的最小元素且 S 是最大元素的情形；最后这种情形利用 e) 处理）。

\begin{enumerate}
\def\labelenumi{\arabic{enumi})}
\setcounter{enumi}{14}
\tightlist
\item
  设 \(\mathbf{T}\) 是一个具有 \(k\) 个元素的有限集。我们用 \(\mathbf{D}_k\) 表示所有元素 \((\alpha_1, \ldots, \alpha_k)\)（属于 \(\mathbb{N}^k\)，满足 \(\alpha_1 + \ldots + \alpha_k = k\) 且 \(\alpha_1 \geqslant \alpha_2 \geqslant \ldots \geqslant \alpha_k\)）组成的集合。对 T 中的每个等价关系 R，我们按如下方式关联元素 \(\lambda(R)\) （属于 \(D_k\)）：若 \(U_1, \ldots, U_l\) 是 R 下的等价类，编号使得 \(\operatorname{Card}(U_1) \geqslant \ldots \geqslant \operatorname{Card}(U_l)\)，令 \(\lambda(R) = (\alpha_1, \ldots, \alpha_k)\)，其中 \(\alpha_i = \operatorname{Card}(U_i)\) （若 \(1 \leqslant i \leqslant l\)），而 \(\alpha_i = 0\) （若 \(l < i \leqslant k\)）。设 \(X_n (n \in N)\) 为不定元；对每个 \(R \in \mathcal{R}(T)\)，令 \(r(R) = \sum_{f} \prod_{t \in T} X_{f(t)}\)，其中求和遍及映射 \(f: T \to N\) （使得 R 是与 \(f\) 相伴的等价关系）组成的集合。令 \(s(R) = \sum_{S \geqslant R} r(S)\) 。
\end{enumerate}

\begin{enumerate}
\def\labelenumi{\alph{enumi})}
\tightlist
\item
  证明 \(r(\mathbf{R}) = \prod_{i=1}^{k} r_i \cdot \mathbf{M}(\lambda(\mathbf{R}))\)，如果有 \(r_i\) 个 \(\mathbf{R}\) 下的等价类的
\end{enumerate}

基数为 \(i\)，对 \(1 \leqslant i \leqslant k\)。（\(\mathbf{M}(\alpha)\) 的定义参见 IV 第 65 页公式 (11)。）

\begin{enumerate}
\def\labelenumi{\alph{enumi})}
\setcounter{enumi}{1}
\item
  证明关系式 \(S(R) = \prod_{i=1}^{k} p_i^{r_i}\)，其中 \(r_i\) 如 a) 中所定义，而 \(p_m = \sum_{n \in N} X_n^m\) 对每个整数 \(m \geqslant 1\) 。
\item
  证明 \(r(\mathbf{R}) = \sum_{\mathbf{S} \geqslant \mathbf{R}} \mu(\mathbf{R}, \mathbf{S}) s(\mathbf{S})\)，其中 \(\mu(\mathbf{R}, \mathbf{S})\) 如练习 14 中那样定义。
\item
  由此推出把 \(s_k\) 表示为 \(p_1, \ldots, p_k\) 的系数为有理数的多项式的一个表达式，其中 \(s_k = \sum_{i_1 \leq \ldots \leq i_k} X_{i_1} \ldots X_{i_k}\) 。
\item
  对 \(\alpha \in \mathbf{D}_k\)，令 \(S'(\alpha) = \prod_{i=1}^{k} s_{\alpha_i}\)。证明 \(S'(\alpha) = \sum_{\beta \in D_k} c_{\alpha\beta}' \cdot M(\beta)\)，其中整数 \(c_{\alpha\beta}'\) 定义如下：设有 \(r_i\)（相应地 \(s_i\)）个 \(\alpha\)（相应地 \(\beta\)）的分量等于 \(i\)（对 \(1 \leqslant i \leqslant k\)）。则 \(c_{\alpha\beta}'\) 是乘积 \(Nr_1!... r_k!s_1!... s_k!/k!\)，其中 N 是 T 上满足下列条件的等价关系对 (R, S) 的数目：\(\lambda(R) = \alpha, \lambda(S) = \beta\)，且 R 下的一个等价类与 S 下的一个等价类的交集至多含有一个元素。由此推出对称关系 \(c_{\alpha\beta}' = c_{\beta\alpha}'\) 。
\item
  证明存在一个对合 \(\alpha \mapsto \tilde{\alpha}\)（在 \(D_k\) 中），它由下述事实刻画：关系 \(i \leqslant \alpha_j\) 与 \(j \leqslant \tilde{\alpha}_i\)（对 \(i\) , \(j\) ，范围从 1 到 \(k\)）是等价的。证明 \(S'(\alpha) = S(\tilde{\alpha})\) 且 \(c_{\alpha\beta}' = c_{\tilde{\alpha}\beta}\)。由此得到关系式 \(c_{\alpha\beta} = c_{\tilde{\alpha}\tilde{\beta}}\)（采用 IV 第 66 页的记号）。
\end{enumerate}

